% Generated from chapter-04.md - do not edit.

\chapter{Negativity Bias}%
\index{negativity bias}\nobreak%

If we want to find better thoughts, it helps to understand why negative ones can be so quick to appear and so hard to let go. Our brains have a natural tendency toward the negative, and knowing why can make it easier to notice those thoughts and gently shift away from them.

You may know about the negative bias. Rick Hansen talks about our survival instinct, that our brains are looking for danger and protecting us. Rick Hanson says that our brains are Teflon for good things and Velcro for anything that looks threatening.   Psychologists call this tendency \textbf{negativity bias}---our brain's built-in habit of noticing, reacting to, and remembering negative experiences more strongly than positive ones.   Negative thoughts feel heavier, sharper, and more convincing than positive ones. They arrive quickly, cling stubbornly, and often feel like the truth itself.

Positive thoughts, by contrast, can seem light, fragile, and easy to lose.\\
\emph{(They don't always know how to make a grand entrance.)}

This is not a personal failure.

It is the human brain at work helping you survive. Thank your brain and think up something positive.

\textbf{Negativity Bias Shows Up in Everyday Life}\index{negativity bias}

The way it works is that when we decide we're going to go to that party we've been invited to, our brains not only think of the good things about the party, but they also worry about many things, like what to wear, and will we be accepted, and will it be any fun at all or is there something else planned for that time. These thoughts come so quickly that we're not even aware of all of them or forget some of them instantly. When we have one thought, other related  ideas come quickly to join them. I believe this is how premonitions are born. If we have a passel of negative thoughts about the party and some undesirable thing happens, we may remember that we thought that might occur! And we say, ``I had a premonition that that might happen.'' This only fuels the false belief that we have power to predict.

\textbf{We play ``Watch Out'' with our friends}\index{watch out}

I would like to add that the negative bias leads us to more than protecting ourselves. We also want to protect our loved ones and friends. When they begin talking about some new venture or an exciting event, we unknowingly scan for the danger in that and warn them to be cautious. This does not go over well. It feels like raining on the parade, a critical comment, when it's just the negative bias. When you have some good news, you would like others to respond enthusiastically about your good fortune in a way that makes you feel happy that you told them. I love it when someone makes a big deal about some delightful occurrence in my life. And yet they seek to save me because their own brain thinks of possible problems.

Negativity Bias appears as worry and pessimism.\index{pessimism}

It appears as remembering one criticism more vividly than ten compliments.

It appears as replaying mistakes long after the moment has passed.

It appears as a host of awful possibilities --- but not real probabilities --- a big difference.

At night, the brain keeps watch.

Many dreams feature being late, lost, unprepared, or failing in some way.

We may wake up uneasy before a single conscious thought forms---the emotional residue still present. My morning thoughts are often grim and I wonder about those vaguely unsettling dreams.

\textbf{Pause for a Moment}

What has your mind been noticing today?

What's been sticking?

A worry?\\
A mistake?\\
Something that might go wrong?

Now gently ask:

\textbf{What's something good that also happened---even if it didn't stick as easily?}

\emph{You may need to hold onto it for a moment.}\\
\emph{(Teflon can be slippery.)}

\section{Thoughts Have Companions}%
\index{companion thoughts}\nobreak%

I am beginning to understand something about thoughts: \textbf{we rarely have just one thought at a time.}

A thought tends to attract other thoughts. One idea connects with another, and then another. Before long, what began as a single thought has gathered a little company around it.

This can be wonderful.

When we are having a happy conversation with someone, one delightful thought may lead to another. We remember something funny. That reminds us of something else we enjoyed. Someone adds another idea. We become curious, playful, interested, appreciative. The conversation begins to sparkle.

One good thought has invited other good thoughts to join it.

We can help that happen. We can stay with the pleasure of the conversation instead of correcting it, criticizing it, making fun of it, or immediately pointing out what might be wrong. We can resist the temptation to say, ``Yes, but . . . '' and instead help the good thought grow.

We can elaborate on delight.

And perhaps that is one of the ways happiness is created---not by having one magnificent happy thought, but by allowing a \textbf{myriad} of pleasant thoughts to gather around something good.\index{myriad of thoughts}

Unfortunately, the same thing happens with negative thoughts.

I may begin with something quite small:

\emph{I'm tired today.}

But that thought may not stay alone.

\emph{I'm too tired to get everything done.}

Then:

\emph{I'm falling behind.}

\emph{Other people manage better than I do.}

\emph{Why can't I handle this?}

\emph{Maybe I can't cope anymore.}

\emph{I'm clumsy.}

\emph{I'm not doing enough.}

\emph{I'm alone in this miserable mess.}

Suddenly, I am no longer dealing with the original thought, \emph{I'm tired today.} I am dealing with a whole collection of thoughts that have attached themselves to it.

The first thought has acquired companions.

Our brains are wonderful connecting machines. One idea brings up another memory, prediction, fear, judgment, comparison, or imagined consequence. The connections can happen so quickly that we may not even notice them. We simply discover that our mood has changed and we are going downhill.

That is why it may be misleading to say, ``Oh, I just had one bad thought.''

Perhaps we almost never have \emph{just one}.

A negative thought can bloom into a \textbf{myriad of related ideas}, each adding another little piece of worry, discouragement, shame, fear, or unhappiness.

Recognizing this gives us another place to intervene.

We can notice the first thought and then become curious about its companions.

\textbf{What have I added to this thought?}

\textbf{What conclusions came rushing in behind it?}

\textbf{What am I predicting?}

\textbf{What am I saying about myself now that wasn't actually contained in the original thought?}

And then, of course, comes the question at the heart of this book:

\textbf{What's a better thought?}

Not necessarily a falsely cheerful thought. Not a thought that denies what is happening. Simply a thought that does not invite an entire unhappy crowd to come marching in behind it.

Perhaps:

\emph{Yes, I'm tired today. That means I need a slower day.}

And then another thought might join that one:

\emph{I know how to rest.}

And another:

\emph{I don't have to do everything.}

And another:

\emph{I can choose one thing that matters and let the rest wait.}

Now I have changed the company my thought is keeping.

This may be something worth remembering: \textbf{thoughts have companions.}

When the companions are delightful, welcome them. Elaborate. Enjoy them. Let one happy thought lead to another and another.

And when the companions are taking you down a dark road, notice them. They are not necessarily truths. They may simply be associations---thoughts joining thoughts because that is what brains do.

Then we have a choice.

We can ask which thoughts we want to invite to stay.

\textbf{NEGATIVITY BIAS SOUNDS LIKE AN ACCUSATION}

One more admonition about negative bias. When I tell people that we all have a negative bias often the person doesn't think of `we all', they believe I'm saying that \emph{they} are negative. Somehow, as soon as you use the word negative the person gets stuck on that and it becomes an insult. They get defensive.

I think we need to rename the negative bias. Give it a new name, something like SAFETY TILT or PROTECTIVE INSTINCT or something that makes it sound like you are doing a positive thing for the sake of your health and well-being when being concerned about unintended consequences. The survival instinct called Safety Tilt makes you a discerning, thoughtful, careful person who makes wonderful decisions. If I have built-in Protective Instinct I become a model for avoiding problems and potential negative consequences. Doesn't that sound much better than Negative Bias?

\textbf{My mind naturally leans toward noticing what could threaten me. I can appreciate that---and still choose where to place my attention.}

If I say:

\textbf{My Survival Scanner is working overtime,}

I may even smile.

If I say:

\textbf{My Watchful Mind has taken over,}

I can recognize what is happening without condemning myself.

If I say:

\textbf{My Protective Instinct is trying to help,}

I can respond with understanding rather than shame.

The name matters because awareness works better when we are not busy defending ourselves.

\section{A Better Relationship With Our Own Minds}

Perhaps this is the real goal.

Not to conquer the negative brain.

Not to become a person who never worries.

Not to force positive thinking.

Not to criticize ourselves for being human.

But to understand:

\textbf{My brain has ancient protective equipment.}

\textbf{It naturally notices danger and difficulty.}

\textbf{It can sometimes overestimate what needs protection.}

\textbf{One negative thought can recruit others.}

\textbf{The Safety Tilt can become a slide.}\index{safety tilt}

And when I notice that happening, I have a choice.

I can ask:

\textbf{Is my Safety Tilt helping me right now?}

\textbf{Has my Survival Scanner found a real problem---or is it simply continuing to scan?}

\textbf{Is my Watchful Mind giving me useful information---or frightening me?}

\textbf{Is there one reasonable action I need to take?}

And then, perhaps, the most important question:

\textbf{What's a better thought?}

Not because the protective mind is bad.

Not because danger is unreal.

Not because we should ignore problems.

But because the ancient system that helped us survive does not always know when its work is done.

We can thank it for trying to protect us.

And then we can look again.

Toward love.

Toward goodness.

Toward what is actually happening now.

Toward the rest of life that the Safety Tilt may have temporarily tilted out of view.
